\section*{\SGAInlineRefTarget{sga3:x:section:7}7. Classification of twisted constant preschemes and of groups of
multiplicative type of finite type in terms of the enlarged fundamental group}
\addcontentsline{toc}{section}{7. Classification in terms of the enlarged
fundamental group}

\setcounter{footnote}{57}

\subsection*{\SGAInlineRefTarget{sga3:x:subsection:7-0}7.0.}

Let \(S\) be a prescheme, which we continue to assume locally noetherian, in
order to ensure that \(S\), and certain preschemes over \(S\) that we shall
consider (notably those étale over \(S\), and more generally those locally of
finite type over \(S\)), are locally connected.

\par\medskip\noindent\SGAInlineRefTarget{sga3:x:proposition:7-0-1}\textbf{Proposition 7.0.1.---}\enspace
\emph{Every twisted constant prescheme \(X\) over \(S\) that is locally
trivial for the \((\mathrm{fppf})\) topology (that is, that is split by a
faithfully flat morphism \(S'\to S\) locally of finite presentation) is
quasi-isotrivial (that is, one may even choose \(S'\to S\) étale and
surjective).}

Indeed, we may suppose that \(S\) is connected, and hence that \(X\) is of
type \(I\), where \(I\) is a fixed set. Thus
\[
  X'=X\times_S S'
\]
is isomorphic to \(I_{S'}\), so \(I_{S'}\) is endowed with descent data
relative to \(S'\to S\); in other words, there is an isomorphism
\[
  I_{S''}\xrightarrow{\sim}I_{S''}
\]
satisfying the usual transitivity condition. Now
\(S''=S'\times_S S'\) is locally noetherian and hence locally connected.
It follows that the automorphisms of \(I_{S''}\) correspond to the sections
of \(G_{S''}\), where \(G=\operatorname{Aut}(I)\) is the group of
permutations of \(I\).

In this way one obtains descent data on \(G_{S'}\) (regarded as a trivial
Galois principal bundle) relative to \(S'\to S\). By \SGARefLink{sga3:x:lemma:5-4}{5.4} these descent data
are effective, and therefore give a Galois principal bundle \(P\) over \(S\)
with group \(G\). By construction, \(P\) represents the functor
\(\operatorname{Isom}_S(I_S,X)\) in the category of preschemes over \(S\)
that are locally noetherian. Consequently, the étale surjective base change
\(P\to S\) splits \(X\), so \(X\) is indeed quasi-isotrivial.

\par\medskip\noindent\SGAInlineRefTarget{sga3:x:remark:7-0-2}\textbf{Remark 7.0.2.---}\enspace
One should bear in mind that, even when \(S\) is the spectrum of a field, it
is not true in general that every twisted constant bundle over \(S\) is
quasi-isotrivial. For example, one may take \(X\) to be the disjoint union
of a sequence of schemes of the form \(\operatorname{Spec}(k_i)\), where
the \(k_i\) are separable extensions of \(k\) of strictly increasing
degrees.

The proof just given shows at the same time that the classification of
quasi-isotrivial twisted constant bundles \(X\) over \(S\) of type \(I\) is
equivalent to that of Galois principal bundles over \(S\) with group
\(G=\operatorname{Aut}(I)\). Indeed, this is an equivalence of categories.

It can be put into a more convenient form, as in SGA~1, V. To this end,
suppose that \(S\) is connected and equipped with a geometric point
\(\xi\). The enlarged fundamental pro-group
\(\Pi=\Pi_1(S,\xi)\) is then defined. For every quasi-isotrivial twisted
constant bundle \(X\) over \(S\), let \(I=X(\xi)\) be its set-theoretic
fiber at \(\xi\). Thus \(X\) is of type \(I\), and is consequently
associated, as above, with a Galois principal bundle
\[
  P=\operatorname{Isom}_S(I_S,X)
\]
over \(S\), with group \(G=\operatorname{Aut}(I)\).

By the definition of \(\Pi\), one therefore obtains a canonical
homomorphism from \(\Pi\) to \(G\), that is, from one of the \(\Pi_i\) to
\(G\). Since \(G\) is the group of permutations of \(I=X(\xi)\), this says
that \(\Pi\) ``acts continuously on \(I=X(\xi)\)'', meaning that the
\(\pi_i\), for \(i\) sufficiently large, act on \(I\) compatibly with the
transition morphisms.

We leave it to the reader to verify that every \(S\)-morphism \(X\to Y\)
between quasi-isotrivial twisted constant bundles over \(S\) induces a map
\(X(\xi)\to Y(\xi)\) compatible with the actions of \(\Pi\), and that the
resulting functor is an equivalence of categories:

\par\medskip\noindent\SGAInlineRefTarget{sga3:x:proposition:7-0-3}\textbf{Proposition 7.0.3.---}\enspace
\emph{Let \(S\) be a locally noetherian connected prescheme, let \(\xi\) be
a geometric point of \(S\), and let
\(\Pi=\Pi_1(S,\xi)\) be the enlarged fundamental pro-group of \(S\) at
\(\xi\). Then the functor}
\[
  X\longmapsto X(\xi)
\]
\emph{is an equivalence between the category of quasi-isotrivial twisted
constant bundles over \(S\) and the category of sets on which \(\Pi\) acts
continuously.}

This functor is compatible with finite sums and finite inverse limits. It
follows, for example, that quasi-isotrivial twisted constant groups (or
rings, and so forth) over \(S\) correspond to ordinary groups
(respectively rings, and so forth) on which the pro-group \(\Pi\) acts
continuously. In particular:

\par\medskip\noindent\SGAInlineRefTarget{sga3:x:corollary:7-0-4}\textbf{Corollary 7.0.4.---}\enspace
\emph{The category of quasi-isotrivial twisted constant commutative groups
over \(S\) is equivalent to the category of ``\(\Pi\)-modules'', that is,
commutative groups \(M\) on which \(\Pi\) acts continuously.}

Using \SGARefLink{sga3:x:corollary:5-7}{5.7}, we now obtain:

\par\medskip\noindent\SGAInlineRefTarget{sga3:x:theorem:7-1}\textbf{Theorem 7.1.---}\enspace
\emph{Let \(S\) be a locally noetherian connected prescheme, let \(\xi\) be
a geometric point of \(S\), and let
\(\Pi=\Pi_1(S,\xi)\) be the enlarged fundamental pro-group of \(S\) at
\(\xi\). Then the functor}
\[
  G\longmapsto
  \operatorname{Hom}_{\kappa(\xi)\text{-}\mathrm{gr}}
    (G_\xi,\mathbf{G}_{m,\xi})
\]
\emph{induces an anti-equivalence between the category of quasi-isotrivial
groups of multiplicative type over \(S\) and the category of
\(\Pi\)-modules.}

Using \SGARefLink{sga3:x:corollary:4-5}{4.5}, we obtain:

\par\medskip\noindent\SGAInlineRefTarget{sga3:x:corollary:7-2}\textbf{Corollary 7.2.---}\enspace
\emph{The preceding functor induces an anti-equivalence between the
category of groups of multiplicative type and of finite type over \(S\),
and the category of \(\Pi\)-modules that are of finite type over
\(\mathbb{Z}\).}

\par\medskip\noindent\SGAInlineRefTarget{sga3:x:example:7-3}\textbf{Example 7.3.---}\enspace
Take for \(S\), for example, a complete rational curve over an
algebraically closed field, having exactly one multiple point with
\(n+1\) distinct branches. By \SGARefLink{sga3:x:example:6-4}{6.4}, the enlarged fundamental group
\(\Pi(S,\xi)\) is a free group on \(n\) generators. Thus, by \SGARefLink{sga3:x:corollary:7-2}{7.2}, the
classification of tori of relative dimension \(m\) over \(S\) is
equivalent to the classification of systems of \(n\) endomorphisms
\(A_1,\ldots,A_n\) of the \(\mathbb{Z}\)-module \(\mathbb{Z}^m\), up to
automorphism of \(\mathbb{Z}^m\). Except when \(n\leq 1\) or \(m\leq 1\),
an explicit classification of such systems seems hopeless. One can at
least define a multitude of nontrivial invariants for such a system, such
as the characteristic polynomials of the \(A_i\).\footnote{\textit{Editors'
note.} In particular, this shows that the tori of relative dimension \(2\)
considered in \SGARefLink{sga3:x:paragraph:1:6}{1.6} are non-isotrivial.}

\par\medskip\noindent\SGAInlineRefTarget{sga3:x:remark:7-4}\textbf{Remark 7.4.---}\enspace
Without any hypothesis on \(S\), it remains true that, for a given ordinary
commutative group \(M\) of finite type over \(\mathbb{Z}\), the category of
groups of multiplicative type of type \(M\) over \(S\) is anti-equivalent
to the category of Galois principal bundles over \(S\) with group
\(G=\operatorname{Aut}_{\mathrm{gr}}(M)\). This follows easily from \SGARefLink{sga3:x:corollary:5-9}{5.9}
and \SGARefLink{sga3:x:corollary:5-10}{5.10}.

\par\medskip\noindent\SGAInlineRefTarget{sga3:x:remark:7-5}\textbf{Remark 7.5.---}\enspace
The theory of the fundamental pro-group sketched in the present two
sections is more advantageously formulated in the setting of general
sites. In that form it applies equally well, for example, to ordinary
topological spaces, and gives a satisfactory theory at least for a locally
connected topological space (not necessarily locally simply connected).
Here too it seems that one cannot make do with a fundamental group, and
that a pro-group is needed. Finally, once one has at one's disposal the
language of topologies and descent (which really lies at the heart of these
questions), the account sketched here is also technically simpler than
that of SGA~1, V, and should therefore in principle replace it.

\section*{\SGAInlineRefTarget{sga3:x:section:8}8. Appendix: Elimination of certain affineness hypotheses}
\addcontentsline{toc}{section}{8. Appendix: Elimination of certain affineness
hypotheses}

Our aim is to prove the following generalization of 4.9.

\par\medskip\noindent\SGAInlineRefTarget{sga3:x:theorem:8-1}\textbf{Theorem 8.1.---}\enspace
\emph{Let \(S\) be a prescheme, and let \(H\) be a flat \(S\)-prescheme in
groups of finite presentation, with connected affine
fibers.\footnote{\textit{Editors' note.} Observe that the hypotheses imply
that \(H\) is separated over \(S\), by VI\textsubscript{B}, \SGARefLink{thm:sga3-VIB-5.3}{Theorem~5.3} or
\SGARefLink{sga3:x:corollary:5-5}{Corollary~5.5}.} Let \(s\in S\), and suppose that \(H_s\) is a torus. Then
there exists an open neighborhood \(U\) of \(s\) such that \(H|_U\) is a
torus.}

We immediately obtain:

\par\medskip\noindent\SGAInlineRefTarget{sga3:x:corollary:8-2}\textbf{Corollary 8.2.---}\enspace
\emph{Let \(H\) be an \(S\)-prescheme in groups, flat and of finite
presentation over \(S\). Then \(H\) is a torus if and only if its fibers
are tori.}

\par\medskip\noindent\SGAInlineRefTarget{sga3:x:remark:8-3}\textbf{Remark 8.3.---}\enspace
Even when \(S\) is the spectrum of a discrete valuation ring, the
hypothesis in 8.1 that the fibers of \(H\) (here, the generic fiber) be
affine cannot be dropped: there are examples of smooth groups over \(S\)
whose generic fiber is an elliptic curve and whose special fiber is
\(\mathbf{G}_m\).

\medskip
\noindent\emph{Proof of 8.1.}
We may clearly suppose that \(S=\operatorname{Spec}(A)\) is affine. By the
standard procedure (cf.\ EGA~IV\textsubscript{2}, 8.8.2), this reduces us
to the case in which \(S\) is moreover noetherian. We begin by proving 8.2
in this case.

By 4.9, we are reduced to proving that \(H\) is affine over \(S\). We may
therefore suppose\footnote{\textit{Editors' note.} By
EGA~IV\textsubscript{2}, 8.8.2 once again.} that \(A\) is noetherian and
local. Since its completion \(\widehat A\) is faithfully flat over \(A\),
descent reduces us to the case in which \(A\) is a complete noetherian
local ring. Let \(S'\) be the normalization of \(S_{\mathrm{red}}\).
Nagata's theorem says that it is finite over \(S\)
(EGA~\(0_{\mathrm{IV}}\), 23.1.5). Moreover, \(S'\to S\) is surjective,
so
\[
  H'=H\times_S S'\longrightarrow H
\]
is finite and surjective. Hence, to prove that \(H\) is affine, it suffices
to show that \(H'\) is affine (EGA~II, 6.7.1).

Replacing \(S\) by a connected component of \(S'\), we are reduced to the
case in which \(A\) is a noetherian (semilocal), normal, integral
ring.\footnote{\textit{Editors' note.} The original has been expanded in
what follows.} Moreover, after possibly replacing \(A\) by its
normalization in a finite separable extension of its field of fractions,
we may suppose that the generic fiber \(H_\eta\) of \(H\) is
diagonalizable, that is, that there is an isomorphism
\[
  u_\eta:H_\eta\xrightarrow{\sim}T_\eta,
  \qquad T=\mathbf{G}_{m,S}^{\,r}.
\]
We now have:

\par\medskip\noindent\SGAInlineRefTarget{sga3:x:lemma:8-4}\textbf{Lemma 8.4.---}\enspace
\emph{Let \(S\) be a locally noetherian normal irreducible prescheme with
generic point \(\eta\). Let \(H\) be a smooth prescheme in groups over
\(S\), with connected fibers, let \(T\) be a prescheme in groups of
multiplicative type and of finite type over \(S\), and let
\(u_\eta:H_\eta\to T_\eta\) be a homomorphism of algebraic groups over
\(\kappa(\eta)\). Then \(u_\eta\) extends to a homomorphism of groups
\(u:H\to T\).}

After possibly replacing \(S\) by its normalization in a finite separable
extension of its function field, we may suppose that \(T\) is
diagonalizable (which is already the case in the application we have in
mind). Then \(T\) is a closed subgroup of a group of the form
\(\mathbf{G}_{m,S}^{\,r}\), reducing us to the case
\(T=\mathbf{G}_{m,S}\).

It remains only to prove that \(u_\eta\), regarded as a rational map from
\(H\) to \(T=\mathbf{G}_{m,S}\), is everywhere defined (for then the
morphism \(u:H\to T\) extending it is necessarily a group homomorphism).
We may regard \(u_\eta\) as an invertible section \(f\) of the structure
sheaf of \(H_\eta\), and must show that it extends to an invertible section
of the structure sheaf of \(H\). Since \(H\) is smooth over the normal
prescheme \(S\), it is normal (SGA~1, II 3.1). It therefore suffices to
find a closed subset \(Z\) of \(H\), of codimension at least \(2\), such
that \(f\) extends to an invertible section of the structure sheaf of
\(H-Z\). By localizing at the codimension-one points of \(S\), this at
once reduces us to the case in which \(S\) is the spectrum of a discrete
valuation ring \(A\).

Let \(t\) be a uniformizer of \(A\), and let \(t'\) be the section of
\(\mathcal O_H\) that it defines, so that the special fiber \(H_0\) equals
\(V(t')\). By hypothesis, \(H_0\) is smooth and connected over the residue
field \(k\). The function \(f\) is then a rational function on \(H\) with
neither zeros nor poles in \(H-H_0\). Since \(H_0=V(t')\) is an
irreducible divisor, there is an integer \(n\in\mathbb{Z}\) such that
\[
  t'^{\,n}f=t^n f
\]
has neither zeros nor poles, and hence is an invertible section of
\(\mathcal O_H\). It therefore defines a morphism
\(v:H\to\mathbf{G}_{m,S}\). Since
\(v_\eta=t'^{\,n}u_\eta\), and \(u_\eta\) is a group homomorphism, \(v\)
sends the identity section of \(H\) to a section of
\(\mathbf{G}_{m,S}\) whose value at the generic point of \(S\) is \(t^n\).
Because this is a section of \(\mathbf{G}_{m,S}\), \(t^n\) must be a unit;
thus \(n=0\). Hence \(v\) extends \(u_\eta\), completing the proof of \SGARefLink{sga3:x:lemma:8-4}{8.4}.

Applying this lemma in the present situation, we obtain a group
homomorphism
\[
  u:H\longrightarrow T=\mathbf{G}_{m,S}^{\,r}
\]
that induces an isomorphism on the generic fibers. We shall prove that
\(u\) is an isomorphism.

\par\medskip\noindent\SGAInlineRefTarget{sga3:x:lemma:8-5}\textbf{Lemma 8.5.---}\enspace
\emph{For every integer \(n>0\) prime to the residual characteristic of
\(A\), the group
\[
  {}_nH=\ker(n\cdot\operatorname{id}_H)
\]
is finite over \(S\).}

If \(n\) is prime to the residual characteristic of \(A\), then it is
prime to all the residual characteristics of the points of \(S\).
Consequently, \(n\cdot\operatorname{id}_H\) induces an étale morphism on
every fiber of \(H\), and hence is itself étale (SGA~1, I 5.9); its kernel
\({}_nH\) is therefore étale over \(S\). On the other hand, \({}_nH\) is
separated over \(S\), because \(H\) is.
\begingroup
\renewcommand{\thefootnote}{\(\ast\)}
\footnote{By Raynaud's theorem, VI\textsubscript{B}, 5.3.}
\endgroup
\addtocounter{footnote}{-1}%
\footnote{\textit{Editors' note.} This point should be expanded in
connection with the modifications in VI\textsubscript{B}, \S5.}
Moreover, all its fibers have the same rank \(n^r\), since the fibers of
\(H\) are tori, all of the same dimension \(r\) (as \(H\) is smooth over
\(S\)). We conclude that \({}_nH\) is finite over \(S\) (SGA~1, I 10.9,
or EGA~IV\textsubscript{4}, 18.2.9).

By the preceding lemma, \(u({}_nH)\) is therefore a closed subset of
\(T\). On the generic fiber \(T_\eta\), this subset is identical to
\({}_n(T_\eta)\), so it contains the closure of that subset, namely
\({}_nT\). For every \(s\in S\), the groups \({}_n(T_s)\), as \(n\)
varies, are dense in \(T_s\); since \(u_s(H_s)\) is a closed subset
(VI\textsubscript{B}, 1.4.2) containing them, we have
\(u_s(H_s)=T_s\). Thus \(u\) is surjective. A surjective homomorphism
between tori of the same dimension over a field is
flat,\footnote{\textit{Editors' note.} This follows, for example, from
VIII 3.2(a). More generally, if \(f:G\to H\) is a surjective morphism of
algebraic groups over a field \(k\), and \(H\) is reduced, then \(f\) is
flat (cf.\ VI\textsubscript{B}, 1.3).} and hence \(u\) induces a flat
morphism on every fiber. Therefore \(u\) is flat (SGA~1, I 5.9).
Consequently \(K=\ker(u)\) is flat over
\(S\),\footnote{\textit{Editors' note.} The original has been expanded in
what follows.} and hence equals the closure of its generic fiber
\(K_\eta\). By construction, \(K_\eta\) is the identity group. Since
\(K\) is separated over \(S\) (because \(H\) is), its identity section is
closed, and it follows that \(K\) is the identity group. Thus \(u\) is a
monomorphism; since we have seen that it is faithfully flat, it is an
isomorphism (cf.\ SGA~1, I 5.1, or EGA~IV\textsubscript{4}, 17.9.1).
This proves that \(H\) is a torus, and completes the proof of
8.2.\footnote{\textit{Editors' note.} At least in the case considered up
to this point, namely \(S\) locally noetherian.}

\par\medskip\noindent\SGAInlineRefTarget{sga3:x:remark:8-5-1}\textbf{Remark 8.5.1.---}\enspace
Instead of invoking 8.5, one can also invoke Zariski's ``Main Theorem'',
which directly implies that \(u\) is an open immersion and hence an
isomorphism.\footnote{\textit{Editors' note.} The editors did not
understand this remark, since they did not see why \(u\) would
\emph{a priori} have finite and surjective fibers.}

To prove \SGARefLink{sga3:x:theorem:8-1}{8.1}, the quasi-isotriviality theorem reduces us to the case in
which \(S\) is local and \(s\) is its closed
point,\footnote{\textit{Editors' note.} Expand this reduction.} and it
remains to prove that, under the other hypotheses, \(H\) is then a torus.
By \SGARefLink{sga3:x:corollary:8-2}{8.2}, already proved, we are reduced to showing that the fibers of \(H\)
are tori. We may suppose that \(S\) is the spectrum of a complete
noetherian local ring \(A\). By \SGARefLink{sga3:x:corollary:3:3}{3.3}, for every \(n>0\) there exists a
group of multiplicative type \(T_n\), finite over \(S\), and an
isomorphism
\[
  (T_n)_s\simeq {}_n(H_s),
\]
where \(s\) is the closed point of \(S\). Arguing as in \SGARefLink{sga3:x:lemma:3:1}{3.1} and using the
fact that \(T_n\) is finite over
\(S\),\footnote{\textit{Editors' note.} Expand this point.} we see that
the preceding isomorphism comes from a homomorphism
\(u_n:T_n\to H\), which is moreover uniquely determined. (Pass to the
limit over the artinian quotients of \(A\).)

Moreover, by the uniqueness properties, when \(m\) is a multiple of \(n\),
the homomorphism \(u_n:T_n\to H\) is obtained from
\(u_m:T_m\to H\) by restriction to
\({}_n(T_m)\simeq T_n\). It follows from \SGARefLink{sga3:ix:corollary:6:6}{IX 6.6} that the \(u_n\) are
monomorphisms, so the \(T_n\) are subgroups of \(H\), and for \(m\) a
multiple of \(n\) one has
\[
  T_n={}_n(T_m).
\]

Thus, for every \(t\in S\), the groups \((T_n)_t\) are subgroups of
\(H_t\), of type \((\mathbb{Z}/n\mathbb{Z})^r\), where
\(r=\dim H_s=\dim H_t\), and when \(m\) is a multiple of \(n\) one has
\[
  (T_n)_t={}_n(T_m)_t.
\]
The fact that \(H_t\) is a torus now follows from:

\par\medskip\noindent\SGAInlineRefTarget{sga3:x:lemma:8-6}\textbf{Lemma 8.6.---}\enspace
\emph{Let \(H\) be a smooth affine algebraic group over an algebraically
closed field \(k\), and let \((T_n)_{n>0}\) be a family of subgroups of
multiplicative type of \(H\), such that for every integer \(n>0\), \(T_n\)
is of type \((\mathbb{Z}/n\mathbb{Z})^r\), and for every multiple \(m\) of
\(n\), one has \(T_n={}_n(T_m)\). Under these conditions, \(H\) contains
a torus of dimension at least \(r\) containing the \(T_n\). Consequently,
if \(H\) is connected and has dimension at most \(r\), then \(H\) is a
torus of dimension \(r\).}

This is an exercise on affine algebraic groups, which we shall handle by
repeated reference to \emph{Bible}. We restrict attention to the \(T_n\)
for \(n\) prime to the characteristic. Let \(K\) be the closure of the
union of the \(T_n\) in \(H\), endowed with the induced reduced structure.
Standard arguments then show that \(K\) is a commutative algebraic
subgroup of \(H\). By \emph{Bible}, \S4.5, Theorem~4, \(K\) is isomorphic
to a product \(K_u\times K_s\), with \(K_u\) ``unipotent'' and \(K_s\)
diagonalizable. Every diagonalizable subgroup of \(K\) is contained in
\(K_s\), so the \(T_n\) are subgroups of \(K_s\), and therefore
\(K=K_s\). Write \(K=D(M)\), where \(M\) is an ordinary commutative group
of finite type over \(\mathbb{Z}\). The inclusion \(T_n\subseteq K\)
means that \(M\) has a quotient group isomorphic to
\((\mathbb{Z}/n\mathbb{Z})^r\). Since this holds for every integer \(n\)
prime to the characteristic of \(k\) (powers of a single fixed prime
would suffice), \(M\) has rank at least \(r\). Thus \(K\) contains a torus
\(T\) of dimension \(r\). When \(H\) is connected of dimension \(r\), it
follows that \(T=H\), completing the proof of \SGARefLink{sga3:x:lemma:8-6}{8.6}. \SGARefLink{sga3:x:theorem:8-1}{Theorem~8.1} is now
proved.

\par\medskip\noindent\SGAInlineRefTarget{sga3:x:remark:8-7}\textbf{Remarks 8.7.---}\enspace
Using \SGARefLink{sga3:x:theorem:8-1}{8.1}, it should not be difficult to give analogous generalizations of
\SGARefLink{sga3:x:corollary:4-7}{4.7} and \SGARefLink{sga3:x:corollary:4-8}{4.8}. A more interesting problem would be to study the situation of
\SGARefLink{sga3:x:theorem:8-1}{8.1} after dropping the hypothesis that the fibers of \(H\) be affine. One
can show that there is then an open neighborhood \(U\) of \(s\) such that
\(H|_U\) is commutative and, for every \(t\in U\), the geometric fiber
\(H_t\) is an extension of an abelian variety by a
torus.\footnote{\textit{Editors' note.} Give a reference here?} Of course,
in questions of this kind one may restrict to the case in which \(S\) is
the spectrum of a discrete valuation ring, \(s\) its closed point, and
\(t\) its generic point.

This result can be generalized as follows. For every smooth connected
algebraic group \(G\) over an algebraically closed field \(k\), a
well-known theorem of Chevalley says that \(G\) is, in a unique way, an
extension of an abelian variety \(A\) by a smooth connected affine group
\(V\). Define the \emph{abelian rank} (respectively, \emph{reductive rank},
\emph{nilpotent rank}, and \emph{semisimple dimension}) of \(G\), denoted
\(\rho_{\mathrm{ab}}(G)\) (respectively,
\(\rho_r(G)\), \(\rho_n(G)\), and \(d_s(G)\)), to be the dimension of
\(A\) (respectively, the dimension of the maximal tori of \(G\), the
dimension of the Cartan subgroups\footnote{\textit{Editors' note.} Recall
the definition of ``Cartan subgroup''.} of \(G\), and the dimension of the
quotient of \(G\), or equivalently of \(V\), by its radical); see
\emph{Bible} for all these notions. Also define the \emph{unipotent rank}
by
\[
  \rho_u(G)=\rho_u(V)
  =\rho_n(G)-\rho_r(G)-\rho_{\mathrm{ab}}(G).
\]
When the ground field is not algebraically closed, the same names and
notations designate the corresponding invariants of \(G_{\bar{k}}\),
where \(\bar{k}\) is an algebraic closure of \(k\).

Now let \(G\) be a smooth group scheme over the spectrum \(S\) of a
discrete valuation ring. Let \(\rho_{\mathrm{ab}}\), and so on
(respectively, \(\rho'_{\mathrm{ab}}\), and so on), denote the invariants
associated with the special fiber (respectively, the generic fiber). Then
one has the inequalities
\[
  \left\{
  \begin{aligned}
    \rho_{\mathrm{ab}}&\leq \rho'_{\mathrm{ab}},\\
    \rho_r+\rho_{\mathrm{ab}}&\leq
      \rho'_r+\rho'_{\mathrm{ab}},\\
    d_s&\leq d'_s,
  \end{aligned}
  \right.
  \qquad
  \left\{
  \begin{aligned}
    \rho_n&\geq \rho'_n,\\
    \rho_u&\geq \rho'_u.
  \end{aligned}
  \right.
\]
Equivalently, if \(G\) is smooth of finite type over an arbitrary base
\(S\), the functions
\[
  s\longmapsto \rho_{\mathrm{ab}}(s),\qquad
  s\longmapsto \rho_{\mathrm{ab}}(s)+\rho_r(s),\qquad
  s\longmapsto d_s(s)
\]
are lower semicontinuous, while the functions
\[
  s\longmapsto \rho_n(s),\qquad
  s\longmapsto \rho_u(s)
\]
are upper semicontinuous.\footnote{\textit{Editors' note.} Give a
reference for these results?}

The same results probably remain valid without assuming \(G\) smooth over
\(S\), but assuming only that it is flat and of finite presentation over
\(S\), provided that, for an algebraic group \(G\) over an algebraically
closed field \(k\), \(\rho_{\mathrm{ab}}(G)\), and so on, are understood
to mean the corresponding invariants of \(G_{\mathrm{red}}^0\).

In this Seminar we present some results of this kind when \(G\) is affine
over \(S\), or more generally has affine fibers. In this case we shall
verify the semicontinuity properties for \(\rho_r\) and \(\rho_n\), hence
for \(\rho_u=\rho_n-\rho_r\), and the continuity of \(\rho_s\) in a
neighborhood of a point \(s\) whose fiber is a reductive
group.\footnotemark[71]

When \(G\) is already assumed commutative, \SGARefLink{sga3:x:corollary:8-2}{8.2} can be generalized as
follows.

\par\medskip\noindent\SGAInlineRefTarget{sga3:x:theorem:8-8}\textbf{Theorem 8.8.---}\enspace
\footnote{\textit{Editors' note.} G.~Prasad and J.-K.~Yu generalized this
result (subject to some additional hypotheses) without assuming \(G\)
commutative, and with ``torus'' replaced by ``reductive group'' in
hypothesis~(a) and in the conclusion; cf.\ \([{\rm PY06}]\),
Theorem~6.2.}%
\emph{Let \(G\) be a commutative \(S\)-prescheme in groups, flat and of
finite presentation over \(S\), with affine connected fibers. Let
\(s\in S\), and suppose that:
\begin{enumerate}
\renewcommand{\labelenumi}{\alph{enumi})}
\item if \(\bar{k}\) denotes an algebraic closure of \(\kappa(s)\), then
  \((G_{\bar{k}})_{\mathrm{red}}\) is a torus;
\item there is a generization \(t\) of \(s\) such that \(G_t\) is smooth
  over \(\kappa(t)\).
\end{enumerate}
Under these conditions, there exists an open neighborhood \(U\) of \(s\)
such that \(G|_U\) is a torus.}

\emph{(Observe moreover that, if one assumes only that \(G_t\) is affine,
respectively connected, for every generization \(t\) of \(s\), it follows
easily that the same conclusion holds for \(t\) in an open neighborhood
of \(s\).)}

\medskip
\noindent\emph{Proof of \SGARefLink{sga3:x:theorem:8-8}{8.8}.}
It suffices to prove that \(G_s\) is smooth over \(\kappa(s)\). Indeed,
since \(G\) is flat and of finite presentation over \(S\), it then follows
that \(G\) is smooth over \(S\) above a neighborhood of \(s\)
(cf.\ \SGARefLink{sga3:x:lemma:3:5}{3.5}), and we are then in the situation of \SGARefLink{sga3:x:theorem:8-1}{8.1}.

To prove that \(G_s\) is smooth over \(\kappa(s)\), the usual procedure
reduces us to the case in which \(S\) is affine and noetherian. Choose a
morphism \(S'\to S\), where \(S'\) is the spectrum of a discrete
valuation ring, sending the closed point to \(s\) and the generic point
\(\eta\) to \(t\). We are thereby reduced to the case in which \(S\) is
itself the spectrum of a discrete valuation ring; we may moreover suppose
that it is complete and has algebraically closed residue field, and that
\(s\) and \(t=\eta\) are respectively the closed point and the generic
point of \(S\).

Thus \(G\) is flat, separated, and of finite type over \(S\); the generic
fiber \(G_\eta\) is smooth and connected, and the special fiber \(G_0\)
is such that
\[
  T_0=(G_0)_{\mathrm{red}}
\]
is a torus. Let \(m>0\) be an integer prime to the residual
characteristic at \(s\), and hence also to that at \(\eta\). Then
\(m\cdot\operatorname{id}_G\) is étale on every fiber
(VII\textsubscript{A}, 8.4), and is therefore étale because \(G\) is flat
over \(S\) (SGA~1, I 5.9). Its kernel \({}_mG\) is thus étale over \(S\);
since \(G\) is separated
\begingroup
\renewcommand{\thefootnote}{\(\ast\)}
\footnote{By Raynaud's theorem, VI\textsubscript{B}, 5.3.}
\endgroup
\addtocounter{footnote}{-1}%
\footnote{\textit{Editors' note.} See Editors' note (62) in 8.5.}
and of finite type over \(S\), the same holds for \({}_mG\). Since
\[
  ({}_mG)_0={}_m(T_0),
\]
its degree is \(m^r\), where
\[
  r=\dim T_0=\dim G_0=\dim G.
\]
It follows that the rank of
\(({}_mG)_\eta={}_m(G_\eta)\) is at least \(m^r\), which already proves,
using \SGARefLink{sga3:x:lemma:8-6}{8.6}, that \(G_\eta\) is a torus of dimension \(r\), since the two
fibers of \({}_mG\) have the same rank, and hence, as in \SGARefLink{sga3:x:lemma:8-5}{8.5}, that
\({}_mG\) is finite over \(S\).\footnote{\textit{Editors' note.} Revise
the preceding sentence.}

Since \(S\) is complete with algebraically closed residue field, the
finite étale cover \({}_mG\) splits completely. Thus, through every point
of \({}_mG_0\) there passes a section of \(G\) over \(S\); in particular,
the set of points of \(G_0\) through which there passes a section of \(G\)
over \(S\) is everywhere dense. We now use the following lemma.

\par\medskip\noindent\SGAInlineRefTarget{sga3:x:lemma:8-9}\textbf{Lemma 8.9.---}\enspace
\emph{Let \(S\) be a locally noetherian regular prescheme of dimension
\(1\), and let \(G\) be an \(S\)-prescheme in groups, flat and locally of
finite type, such that \(G_\eta\) is smooth over \(\kappa(\eta)\) for
every maximal point \(\eta\) of \(S\) (which implies that \(G\) is
reduced). Suppose that the normalization \(X\) of \(G\) is finite over
\(G\) (this is the case, by Nagata, when \(S\) is the spectrum of a
complete discrete valuation ring; cf.\ EGA~IV\textsubscript{2}, 7.7.4),
and let \(G'\) be the open subset of \(X\) formed by the points at which
\(X\) is smooth over \(S\). With these notations:
\begin{enumerate}
\renewcommand{\labelenumi}{\alph{enumi})}
\item If the projection \(G'\to S\) is surjective, then \(G'\) carries a
  unique structure of an \(S\)-prescheme in groups for which the
  canonical morphism \(G'\to G\) is a group homomorphism.
\item Suppose that, for every closed point \(s\) of \(S\), the set of
  points of \(G_s^0\) through which an étale quasi-section passes is dense
  in \(G_s^0\) for the Zariski topology. Then \(X\) is regular, the
  hypotheses of \textup{(a)} hold, and the map
  \[
    \Gamma(G'/S)\longrightarrow\Gamma(G/S)
  \]
  is bijective.
\end{enumerate}}

In the case that concerns us, this lemma applies and gives a group
homomorphism \(u:G'\to G\), where \(G'\) is smooth over \(S\) and
\[
  u_\eta:G'_\eta\xrightarrow{\sim}G_\eta
\]
is an isomorphism. After possibly replacing \(G'\) by an open subgroup,
we may suppose that \(G'_0\) is connected. Since
\(u_0:G'_0\to G_0\) induces a surjective morphism
\(G'_0\to T_0\), where \(T_0\) is a torus of the same dimension as
\(G'_0\), it follows easily that \(G'_0\) is a torus (for example, by
using \SGARefLink{sga3:x:lemma:8-6}{8.6}, or by referring to \emph{Bible}, \S7.3, Theorem~3(a)).
Consequently, by \SGARefLink{sga3:x:corollary:8-2}{8.2}, \(G'\) is a torus. By \SGARefLink{sga3:ix:theorem:6:8}{IX 6.8},
\(\ker u\) is then a subgroup of multiplicative type of \(G'\); since its
generic fiber is the identity group, it is itself the identity group, and
therefore \(u\) is a monomorphism. It follows from \SGARefLink{sga3:viii:proposition:7:9}{VIII 7.9} that \(u\) is
an immersion. Since it is surjective and \(G\) is reduced, \(u\) is an
isomorphism, completing the proof of \SGARefLink{sga3:x:theorem:8-8}{8.8}.

It remains to prove \SGARefLink{sga3:x:lemma:8-9}{8.9}. For part \SGARefLink{sga3:viii:proposition:7:9:item:a}{(a)}, the uniqueness of the group law on
\(G'\) for which \(u:G'\to G\) is a group homomorphism is clear, since
the group law on the generic fiber of \(G'\) is known (we may suppose
\(S\) irreducible). For existence, we reduce easily to the case in which
\(S\) is local, the spectrum of a discrete valuation ring \(A\). By
uniqueness, and because integral closure commutes with an étale extension
of the base, we may make étale base extensions on \(S\). This reduces us
to the case in which \(A\) is ``strictly local'', that is, henselian with
separably closed residue field. The same reduction applies to part (b);
under the hypothesis made there, ``étale quasi-section'' can now be
replaced by ``section''.

Since \(G'\) is smooth over \(S\) and \(X\) is normal, the prescheme
\(G'\times_S X\) is normal (being smooth over the normal prescheme \(X\)).
Thus the composite morphism
\[
  G'\times_S X\longrightarrow G\times_S G\longrightarrow G
\]
factors as
\[
  p:G'\times_S X\longrightarrow X.
\]
We shall prove that the restriction of this morphism to the open subset
\(G'\times_S G'\) of \(G'\times_S X\) induces a morphism
\[
  G'\times_S G'\longrightarrow G'.
\]

It is therefore enough to show that \(G'_0\times_k G'_0\) maps into the
open subset \(G'_0\) of \(X_0\). It suffices to check that, for every
\(k\)-valued point \(g'_0\) of \(G'_0\), the morphism
\[
  \tag{\(\times\)}\label{sga3:x:lemma:8-9:equation:times}
  h'_0\longmapsto p_0(g'_0,h'_0)
\]
from \(G'_0\) to \(X_0\) takes its values in \(G'_0\). Since \(G'\) is
smooth over \(S\) and \(S\) is henselian, every such \(g'_0\) is induced
by a section \(g'\) of \(G'\) over \(S\). The morphism \((\times)\) is
then induced by the morphism
\[
  h'\longmapsto p(g',h')
\]
from \(G'\) to \(X\), which itself is obtained by transport of structure
from the automorphism \(h\mapsto g\cdot h\) of \(G\), namely left
translation by the section \(g\) of \(G\) that is the image of \(g'\).
Thus \(h'\mapsto p(g',h')\) is itself an automorphism of \(X\), and hence
maps \(G'\) into \(G'\). This proves the assertion.

It remains to prove that the composition law thereby obtained on \(G'\)
is a group law. Associativity follows at once from associativity on the
generic fiber (which is isomorphic to that of \(G\)). Moreover, the
inversion automorphism of the \(S\)-prescheme \(G\) induces an automorphism
of \(X\), which therefore preserves \(G'\) and induces an automorphism
\(\sigma\) of \(G'\). One verifies that \(\sigma\) has the properties of
an inverse for the composition law on \(G'\). Indeed, this again amounts
to verifying the commutativity of certain diagrams involving fiber
powers of \(G'\) over \(S\); since these are smooth over \(S\), it
suffices to check commutativity on the generic fiber, where it is clear.
This proves part (a) of 8.9.

We prove part (b). Let \(Z'\) be the set of \(x\in X\) for which
\(\mathcal O_{X,x}\) is not regular. It is closed by a theorem of Nagata
(EGA~IV\textsubscript{2}, 6.12.6), and is clearly contained in \(X_0\).
Let \(Z\) be its image in \(G\), which is therefore a closed subset of
\(G_0\). Then \(Z\) is nowhere dense in \(G_0\), that is, it contains no
maximal point \(y\) of \(G_0\). Indeed, \(G_0\) is defined in \(G\) by an
equation \(t=0\), where \(t\) is a uniformizer of the discrete valuation
ring \(A\). By the Hauptidealsatz, \(\mathcal O_{G,y}\) has dimension
\(1\). Hence, for every \(x\in X\) above \(y\),
\(\mathcal O_{X,x}\) has dimension \(1\), and is therefore a discrete
valuation ring, since \(X\) is normal and hence regular in codimension
\(1\).

On the other hand, for every section \(g\) of \(G\) over \(S\), the set
\(Z'\) is plainly stable under the automorphism of \(X\) obtained by
transport of structure from left translation by \(g\) on \(G\). Thus
\(Z\) is stable under the left translation on \(G_0\) defined by \(g_0\).
By hypothesis, the set of these \(g_0\) is dense in \(G_0^0\). Since
\(Z\) is stable under translation by these \(g_0\) and is a nowhere-dense
closed subset, it follows at once that \(Z=\varnothing\), whence
\(Z'=\varnothing\), so \(X\) is regular. Consequently \(X\) is smooth
over \(S\) at every point through which a section passes. Every section
of \(G\) over \(S\) lifts uniquely to a section of \(X\) over \(S\), and
hence to a section of \(G'\) over \(S\). This applies in particular to the
identity section, proving that the image of \(G'\) in \(S\) is all of
\(S\), that is, that the hypothesis of (a) holds. This completes the
proof of \SGARefLink{sga3:x:lemma:8-9}{8.9}, and hence of \SGARefLink{sga3:x:theorem:8-8}{8.8}.

\section*{\SGAInlineRefTarget{sga3:x:section:9}9. Addenda}
\addcontentsline{toc}{section}{9. Addenda}
\setcounter{section}{9}
\setcounter{subsection}{0}
\footnote{\textit{Editors' note.} This additional section was written by
Fabrice Orgogozo (following indications by Ofer Gabber).}

\subsection{Simplicial sets, topoi, groupoids, and topological spaces}

\par\medskip\noindent\SGAInlineRefTarget{sga3:x:notation:9-1-1}\textbf{Notations 9.1.1.---}\enspace
\emph{(1) Let \(E\) be a simplicial set. The following objects can be
associated with it:
\begin{enumerate}
\renewcommand{\labelenumi}{(\arabic{enumi})}
\setcounter{enumi}{1}
\item a topos \(T=E^{\sim}\), obtained from the topoi naturally
  associated with the sets \(E_i\), by the procedure described in
  \([{\rm Del74}]\), 6.3.1 (see also \([{\rm Ill72}]\), VI.5.2 and
  SGA~4, VI.7);
\item a groupoid \(G=\Pi(E)\), whose objects are the elements of \(E_0\)
  (the ``vertices'') and whose arrows are defined in \([{\rm GZ67}]\),
  II.7;
\item a topological space \(X=|E|\), a cellular complex called the
  ``geometric realization'' (or ``topological realization'')
  (cf.\ \emph{loc.\ cit.}, III.1).
\end{enumerate}}

Observe that a sheaf on \(T\) is nothing other than a simplicial set over
\(E\).

\subsection{Locally constant sheaves; descent data}

\par\medskip\noindent\SGAInlineRefTarget{sga3:x:definition:9-2-1}\textbf{Definitions 9.2.1.---}\enspace
A \emph{locally constant sheaf} on
\begin{enumerate}
\renewcommand{\labelenumi}{(\arabic{enumi})}
\item a simplicial set \(E\) is a morphism of simplicial sets
  \(E'\to E\) such that, for every \(i\in\mathbb{N}\) and every
  \(e\in E_i\), the face operators \(d\) induce isomorphisms
  \[
    E'_e\xrightarrow{\sim}E'_{d(e)}
  \]
  between the fibers (cf.\ \([{\rm AM69}]\), \S10);
\item a topos \(T\) is an object \(F\) of \(T\) such that there exist an
  epimorphism \(U\to 1\) and an isomorphism
  \[
    F\times U\simeq f^*I\times U,
  \]
  where \(I\) is a set and \(f:T\to\mathbf{Ens}\) is the final morphism
  (cf.\ SGA~4, IX.2);
\item a groupoid \(G\) is a presheaf on \(G\), that is, a contravariant
  functor from \(G\) to the category of sets
  (cf.\ \([{\rm GZ67}]\), Appendix I.1.2);
\item a topological space \(X\) is a sheaf of sets on \(X\) that is
  locally constant in the usual sense.
\end{enumerate}
Finally:
\begin{enumerate}
\renewcommand{\labelenumi}{(\arabic{enumi})}
\setcounter{enumi}{4}
\item a \emph{descent datum} on a simplicial set \(E\) consists of a
  sheaf \(F\) on \(E_0^{\sim}\) (that is, a set-valued function on
  \(E_0\)) and an isomorphism
  \[
    \alpha:p_1^*F\xrightarrow{\sim}p_2^*F,
  \]
  where \(p_1,p_2:E_1^{\sim}\to E_0^{\sim}\) are the morphisms induced by
  the faces, satisfying the usual cocycle relation
  (cf.\ \SGARefLink{sga3:iv:definition:2:1}{Exposé~IV, 2.1}%
  \begingroup
  \renewcommand{\thefootnote}{1}
  \footnote{PP: I have replaced ``FGA, \emph{Technique de descente} I,
  \SGARefLink{sga3:i:subsection:1-6}{1.6}'' by ``\SGARefLink{sga3:iv:definition:2:1}{Exposé IV, 2.1}''.}
  \endgroup
  \addtocounter{footnote}{-1}%
  \ and below).
\end{enumerate}

The morphisms between these five kinds of objects, as well as the
associated inverse-image functors, are defined in the evident way.

\subsection{Some equivalences of categories}

Let \(E\) be a simplicial set, and let \(T\), \(G\), and \(X\) be,
respectively, the associated topos, groupoid, and topological space.

\par\medskip\noindent\SGAInlineRefTarget{sga3:x:proposition:9-3-1}\textbf{Proposition 9.3.1.---}\enspace
\emph{The categories of locally constant sheaves on \(E\), \(T\), \(G\),
and \(X\), as well as the category of descent data on \(E\), are
equivalent.}

\medskip
\noindent\emph{Sketch of proof.}
Denote by \SGARefLink{sga3:x:proposition:9-3-1:item:1:03}{(1)} through \SGARefLink{sga3:x:proposition:9-3-1:item:5}{(5)} the categories of objects defined in the
preceding paragraph.

\begin{itemize}
\item[\((1)\Longleftrightarrow(5)\).]
  This is a special case of \([{\rm AM69}]\), 10.6 (see also
  \([{\rm GZ67}]\), Appendix \SGARefLink{sga3:x:corollary:2:3}{I.2.3}, \([{\rm Fri82}]\), ?.5.6, or
  \([{\rm Ill72}]\), VI.8.1.6). An equivalence of categories is given by
  the functor that associates with the object \(E'\to E\) the pair
  \((E'_0,\alpha)\), where \(E'_0\) is regarded as a sheaf on \(E_0\)
  and \(\alpha\) is the unique isomorphism
  \[
    p_1^*E'_0\xrightarrow{\sim}p_2^*E'_0
  \]
  whose fiber at each \(y\in E_1\), with images denoted \(x_1\) and
  \(x_2\) under the two projections, is given by the isomorphisms
  \[
    (E'_0)_{x_1}\xleftarrow{\sim}(E'_1)_y
      \xrightarrow{\sim}(E'_0)_{x_2}.
  \]
\item[\(\SGAInlineRefTarget{sga3:x:proposition:9-3-1:item:5}(5)\Longleftrightarrow(3)\).]
  This is evident: one of the two relations defining the morphisms in the
  groupoid \(G\) associated with \(E\) is a cocycle relation.
\item[\((1)\Longleftrightarrow(4)\).]
  Cf.\ \([{\rm GZ67}]\), Appendix I.3.2.1.
\item[\(\SGAInlineRefTarget{sga3:x:proposition:9-3-1:item:1:03}(1)\Longleftrightarrow(2)\).]
  One must show that an object over \(E\) is a locally constant sheaf in
  the simplicial sense if and only if it is one as a sheaf on
  \(T=E^{\sim}\). This follows from the fact that the locally constant
  sheaves \(F_\bullet\) on \(T\) are precisely the cartesian sheaves,
  that is, those for which the arrows
  \[
    ([n]\to[m])^*F_m\longrightarrow F_n
  \]
  are isomorphisms. This last assertion is a special case of a general
  fact about simplicial topoi, together with the fact that every sheaf on
  \(E_0^{\sim}\) is locally constant. (See also \([{\rm Ill72}]\),
  VI.8.1.6.)
\end{itemize}
\hfill Q.E.D.

For every group \(H\), the preceding equivalences of categories induce
equivalences between the categories of \(H\)-torsors. Here these are
locally constant sheaves equipped with an action of \(H\), whose fibers
are isomorphic to \(H\) acting on itself by translation.

\subsection{Fundamental groups and groupoids}

It follows from the preceding equivalences of categories that, for every
group \(H\) and every simplicial set \(E\), the sets of isomorphism
classes of \(H\)-torsors
\[
  H^1(E,H),\qquad
  H^1(E^{\sim},H),\qquad
  \pi_0\bigl(\operatorname{Hom}(\Pi(E),BH)\bigr),\qquad
  H^1(|E|,H)
\]
are naturally in bijection. Recall that \(BH\) denotes the one-object
groupoid associated with \(H\), and that \(\pi_0\) is the functor that
associates with a category the set of isomorphism classes of its objects.

Likewise, if \(e\) is a point of \(E\), the preceding equivalences induce
bijections between the sets of isomorphism classes of \(H\)-torsors
trivialized over \(e\), denoted respectively by
\[
  H^1(E\mathbin{\mathrm{rel}}e,H),\qquad
  H^1(E^{\sim}\mathbin{\mathrm{rel}}e^{\sim},H),\qquad
  \operatorname{Hom}\bigl(\pi_1(\Pi(E),e),H\bigr),\qquad
  H^1(|E|\mathbin{\mathrm{rel}}|e|,H).
\]
Recall that
\[
  \pi_1(\Pi(E),e)=\operatorname{Isom}_{\Pi(E)}(e,e).
\]
For variable \(H\), these functors are represented, in the connected
case, by a group denoted \(\pi_1(E,e)\). The group \(\pi_1(E,e)\) is
isomorphic to \(\pi_1(\Pi(E),e)\) and to \(\pi_1(|E|,|e|)\), and hence
also to the fundamental group of a simplicial set as defined by Kan
(cf., for example, \([{\rm May67}]\), 16.1, or \([{\rm Ill72}]\),
I.2.1.1). Recall also that \(H^1(E,H)\) is in turn isomorphic to the set
\(H^1(\pi_1(E,e),H)\) of homomorphisms to \(H\) modulo conjugation, also
denoted
\[
  \operatorname{Hom\,ext}(\pi_1(E,e),H).
\]

\subsection{Cones}

\par\medskip\noindent\SGAInlineRefTarget{sga3:x:definition:9-5-1}\textbf{Definitions 9.5.1.---}\enspace
\emph{\SGAInlineRefTarget{sga3:x:definition:9-5-1:item:1}(1) Let \(f:E'\to E\) be a morphism of simplicial sets. Recall
(cf., for example, \([{\rm Del74}]\), 6.3.1) that \(C(f)\) denotes the
pointed simplicial set whose underlying set in degree \(n\geq0\) is}
\[
  E_n\amalg\left(\coprod_{i<n}E'_i\right)\amalg\star,
\]
\emph{where \(\star\) is a singleton. We leave it to the reader to define
the simplicial maps; the faces in degrees at most two are recalled below.
(See also \([{\rm GZ67}]\), VI.2, for a pointed variant.) The category of
locally constant sheaves on \(C(f)\) is equivalent to the category of
locally constant sheaves on \(E\) equipped with a trivialization of the
inverse image on \(E'\). The simplicial set \(C(f)\) is naturally pointed
by \(\star\to C(f)\).}

\emph{(2) Let \(f:T'\to T\) be a morphism of topoi. Denote by \(C(f)\) the
topos whose objects are quintuples}
\[
  \bigl(F,F',A,\alpha:f^*F\to F',
                 \beta:A\to\Gamma(T',F')\bigr),
\]
\emph{where \(F\) (respectively \(F'\)) is an object of \(T\)
(respectively \(T'\)), \(A\) is a set, and \(\alpha\) (respectively
\(\beta\)) is a morphism in \(T'\) (respectively in \(\mathbf{Ens}\)).
(See also \([{\rm Del80}]\), 4.3.4 and \([{\rm Ill72}]\), III.4 for a
variant of this construction.) The category of locally constant sheaves
on \(C(f)\) is equivalent to the category of locally constant sheaves on
\(T\) equipped with a trivialization of the inverse image on \(T'\). The
topos \(C(f)\) is naturally pointed by the fiber functor that sends the
quintuple above to the set \(A\).}

\emph{(3) Let \(f:G'\to G\) be a morphism of groupoids. Denote by \(C(f)\)
the colimit of the diagram}
\[
  G\longleftarrow G'\longrightarrow B1,
\]
\emph{where \(B1\) is the one-object category (one object, one arrow).
The category of locally constant sheaves on \(C(f)\) is equivalent to the
category of locally constant sheaves on \(G\) equipped with a
trivialization of the inverse image on \(G'\).}

\emph{(4) Let \(f:X'\to X\) be a morphism of topological spaces. Denote by
\(C(f)\) the colimit of the diagram}
\[
  \star\longleftarrow X'\times\{0\}
  \longrightarrow X'\times[0,1]
  \longleftarrow X'\times\{1\}
  \longrightarrow X.
\]
\emph{The category of locally constant sheaves on \(C(f)\) is equivalent
to the category of locally constant sheaves on \(X\) equipped with a
trivialization of the inverse image on \(X'\). The topological space
\(C(f)\) is naturally pointed by \(\star\to C(f)\).}

\subsubsection*{9.5.2. Descent data on the cone of a simplicial map}

Let \(f:E'\to E\) be a morphism of simplicial sets, and let \(C(f)\) be
its cone (cf.\ above, \SGARefLink{sga3:x:definition:9-5-1:item:1}{(1)}). Use \(p\) (respectively \(q,r\)) to denote the
face maps of \(E'\) (respectively \(E,C(f)\)). Before stating the
proposition below, we describe the faces \(r\) in degrees at most two in
terms of \(p\) and \(q\). By convention, \(p_{ij}\) (respectively \(p_i\))
is the face map
\[
  E'_2=E'_{\{1,2,3\}}\longrightarrow
  E'_1=E'_{\{1,2\}}
  \quad
  \left(\text{respectively }
  E'_1=E'_{\{1,2\}}\longrightarrow E'_0=E'_{\{1\}}\right)
\]
corresponding to the increasing map
\[
  \{1,2\}\longrightarrow\{1,2,3\}
  \quad
  \left(\text{respectively }
  \{1\}\longrightarrow\{1,2\}\right)
\]
whose image is \(\{i,j\}\) (respectively
\(\{i\}\)).%
\begingroup
\renewcommand{\thefootnote}{2}
\footnote{This convention departs from the present simplicial convention
but is closer to the notation used by Grothendieck in his theory of
descent.}
\endgroup
\addtocounter{footnote}{-1}%
The same convention is adopted for the faces of \(E\) and \(C(f)\).

The morphism
\[
  r_1:C(f)_1=E_1\amalg E'_0\amalg\star
  \longrightarrow C(f)_0=E_0\amalg\star
\]
(respectively \(r_2\)) is:
\begin{itemize}
\item \(q_1:E_1\to E_0\) (respectively \(q_2:E_1\to E_0\)) on \(E_1\);
\item \(E'_0\to\star\) (respectively \(f_0:E'_0\to E_0\)) on \(E'_0\);
\item \(\star\to\star\) on \(\star\).
\end{itemize}
Likewise, the morphism
\[
  r_{21}:C(f)_2=E_2\amalg E'_1\amalg E'_0\amalg\star
  \longrightarrow C(f)_1=E_1\amalg E'_0\amalg\star
\]
(respectively \(r_{32},r_{31}\)) is:
\begin{itemize}
\item \(q_{21}:E_2\to E_1\) (respectively \(q_{32},q_{31}\)) on \(E_2\);
\item \(p_1:E'_1\to E'_0\) (respectively
  \(f_1:E'_1\to E_1\), \(p_2:E'_1\to E'_0\)) on \(E'_1\);
\item \(E'_0\to\star\) (respectively
  \(\operatorname{Id}:E'_0\to E'_0\),
  \(\operatorname{Id}:E'_0\to E'_0\)) on \(E'_0\).
\end{itemize}

Let \(H\) be a descent datum on \(C(f)\), that is, a sheaf \(H_0\) on
\[
  C(f)_0=E_0\amalg\star
\]
equipped with an isomorphism
\[
  \gamma:r_1^*H_0\xrightarrow{\sim}r_2^*H_0
\]
between its inverse images on
\[
  C(f)_1=E_1\amalg E'_0\amalg\star,
\]
satisfying the cocycle relation
\[
  r_{31}^*\gamma=r_{32}^*\gamma\circ r_{21}^*\gamma.
\]

The restriction of \(\gamma\) to \(E_1\) (respectively \(E'_0\)) is a
descent datum
\[
  \beta:q_1^*G_0\xrightarrow{\sim}q_2^*G_0,
\]
where \(G_0\) is the restriction of \(H_0\) to \(E_0\) (respectively, a
trivialization
\[
  t:H_0(\star)\xrightarrow{\sim}f^*G_0=:F_0,
\]
where \(H_0(\star)\) is the constant sheaf with stalk \(H_0(\star)\)).
The restriction of the cocycle relation satisfied by \(\gamma\) to
\(E_2\) (respectively \(E'_1\)) is the cocycle relation for \(\beta\)
(respectively,
\[
  p_2^*(t)=f_1^*(\beta)\circ p_1^*(t)).
\]
This latter relation says that the trivialization \(t\) is compatible
with the descent datum \(f_1^*(\beta)\) induced by \(\beta\) on \(F_0\).
It follows that:

\par\medskip\noindent\SGAInlineRefTarget{sga3:x:proposition:9-5-3}\textbf{Proposition 9.5.3.---}\enspace
\emph{Let \(f:E'\to E\) be a morphism of simplicial sets, and let \(C(f)\)
be its cone, pointed by \(\star\). The category of descent data on
\(C(f)\) trivialized over \(\star\) is equivalent to the category of
descent data on \(E\) equipped with a trivialization of the descent datum
induced on \(E'\).}

\subsection{Representability of the functor
\(\overline{\pi}^{\,1}(S'/S,\xi;{-})\)}

The notation is that of page 90, and from now on we suppose that the
connected components of
\[
  S_0=S'
  \qquad\text{and}\qquad
  S_1=S'\times_S S'
\]
are open. Under this hypothesis we have the following explicit
combinatorial description, which follows immediately from the preceding
recall.

\par\medskip\noindent\SGAInlineRefTarget{sga3:x:proposition:9-6-1}\textbf{Proposition 9.6.1.---}\enspace
\emph{For every group \(G\), the set
\(\overline{\pi}^{\,1}(S'/S,\xi;G)\) of isomorphism classes of descent
data equipped with a trivialization over \(\xi\) is naturally in
bijection with the relative nonabelian cohomology set}
\[
  H^1(K\mathbin{\mathrm{rel}}k,G)
  :=H^1(\widetilde K\mathbin{\mathrm{rel}}\star,G).
\]

It follows that this functor is representable by a group, denoted
\(\pi_1(K,k)\), the \emph{relative fundamental group}, as soon as the
simplicial set
\[
  \widetilde K=C(k\to K)
\]
is connected. One verifies immediately that this holds if and only if
the map \(\pi_0(k)\to\pi_0(K)\) is surjective. It is enough, in
particular, that the simplicial set \(K\) be connected. As indicated in
the text, this is the case when the scheme \(S\) is connected and the
morphism \(S'\to S\) is universally submersive.

\subsection{Contractibility of \(k\); counterexample to the injectivity of
\(k_\bullet\to K_\bullet\)}

The following two propositions make Editors' note (53) more precise.

\par\medskip\noindent\SGAInlineRefTarget{sga3:x:proposition:9-7-1}\textbf{Proposition 9.7.1.---}\enspace
\emph{Let \(\kappa\) be an algebraically closed field, and let \(X\) be a
connected \(\kappa\)-scheme. For every integer \(i\geq0\), let \(X_i\)
denote the \((i+1)\)-st fiber power of \(X\) over \(\kappa\). For every
integer \(i\geq0\), the canonical map}
\[
  \pi_0(X_i)\longrightarrow\pi_0(X)^{i+1}
\]
\emph{is a bijection. In particular, the simplicial set
\(\pi_0(X_\bullet)\) is contractible.}

It suffices to prove the following lemma, which follows by passage to the
limit from the Künneth formula. [Expand?]

\par\medskip\noindent\SGAInlineRefTarget{sga3:x:lemma:9-7-2}\textbf{Lemma 9.7.2.---}\enspace
\emph{Let \(\kappa\) be an algebraically closed field, and let \(X\) and
\(Y\) be two connected \(\kappa\)-schemes. Then the fiber product
\(X\times_\kappa Y\) is connected.}

\par\medskip\noindent\SGAInlineRefTarget{sga3:x:proposition:9-7-3}\textbf{Proposition 9.7.3.---}\enspace
\emph{Let \(X\) be a connected scheme, let \(X'\to X\) be a finite étale
morphism, and let \(\bar{x}\) be a geometric point of \(X\) localized at
a point \(x\). The canonical morphism}
\[
  \pi_0(X'_{\bar{x}})\longrightarrow\pi_0(X')
\]
\emph{is surjective.}

Indeed, the image of every connected component of \(X'\) is open and
closed in \(X\), and therefore meets the fiber \(X'_x\). Thus
\[
  \pi_0(X'_x)\longrightarrow\pi_0(X')
\]
is surjective, as is
\[
  \pi_0(X'_{\bar{x}})\longrightarrow\pi_0(X'_x).
\]

\section*{Bibliography}
\addcontentsline{toc}{section}{Bibliography}
\footnote{\textit{Editors' note.} Additional references cited in this
Exposé; the references \([{\rm AM69}]\) and following concern the
Addenda.}

\begin{description}
\item[\({[{\rm BAC}]}\)]
  N.~Bourbaki, \emph{Algèbre commutative}, Chap.~I--IV, Hermann, 1961;
  Masson, 1985; Springer-Verlag, 2006.
\item[\({[{\rm BEns}]}\)]
  N.~Bourbaki, \emph{Théorie des ensembles}, Chap.~I--IV, Hermann, 1970.
\item[\({[{\rm Pes66}]}\)]
  C.~Peskine, \emph{Une généralisation du ``Main Theorem'' de Zariski},
  Bull.\ Sci.\ Math.\ \textbf{90} (1966), 119--127.
\item[\({[{\rm PY06}]}\)]
  G.~Prasad and J.-K.~Yu, \emph{On quasi-reductive group schemes},
  J.\ Alg.\ Geom.\ \textbf{15} (2006), 507--549.
\item[\({[{\rm Ray70}]}\)]
  M.~Raynaud, \emph{Anneaux locaux henséliens}, Lect.\ Notes Maths.\
  \textbf{169}, Springer-Verlag, 1970.
\item[\({[{\rm AM69}]}\)]
  M.~Artin and B.~Mazur, \emph{Etale homotopy}, Lect.\ Notes Maths.\
  \textbf{100}, Springer, 1969.
\item[\({[{\rm Del74}]}\)]
  P.~Deligne, \emph{Théorie de Hodge III}, Publ.\ Math.\ IHÉS
  \textbf{44} (1974), 5--77.
\item[\({[{\rm Del80}]}\)]
  P.~Deligne, \emph{La conjecture de Weil II}, Publ.\ Math.\ IHÉS
  \textbf{52} (1980), 137--252.
\item[\({[{\rm Fri82}]}\)]
  E.~M.~Friedlander, \emph{Etale homotopy of simplicial schemes},
  Princeton Univ.\ Press, 1982.
\item[\({[{\rm GZ67}]}\)]
  P.~Gabriel and M.~Zisman, \emph{Calculus of fractions and homotopy
  theory}, Springer, 1967.
\item[\({[{\rm Ill72}]}\)]
  L.~Illusie, \emph{Complexe cotangent et déformations I \& II},
  Lect.\ Notes Maths.\ \textbf{239} \& \textbf{283}, Springer,
  1971 \& 1972.
\item[\({[{\rm Kan58}]}\)]
  D.~Kan, \emph{A combinatorial definition of homotopy groups},
  Ann.\ of Math.\ \textbf{67} (1958), 282--312.
\item[\({[{\rm May67}]}\)]
  J.~P.~May, \emph{Simplicial objects in algebraic topology},
  Univ.\ of Chicago Press, 1967.
\end{description}
